\section{Model}
\label{sec:model}

There are two jurisdictions, indexed by $i\in\{A,B\}$, and one incumbent manufacturing firm in each jurisdiction. Plant locations are fixed. The model is intended to capture industries in which production is geographically anchored but product-market competition is broader than the jurisdiction itself. Information is complete and all primitives are common knowledge.

\subsection{Timing}

The game has three stages. In Stage 1, local governments simultaneously choose a firm-specific green-investment subsidy $s_i$ and shared green public infrastructure $h_i$. In Stage 2, firms simultaneously choose conventional cost-reducing investment $x_i$ and green investment $g_i$. In Stage 3, firms compete in quantities. We solve for a subgame-perfect Nash equilibrium by backward induction.

\subsection{Product market}

A representative consumer has gross utility
\begin{equation}
V(q_A,q_B)
=a(q_A+q_B)-\frac12(q_A^2+q_B^2)-\theta q_Aq_B,
\label{eq:grossutility}
\end{equation}
which implies inverse demand
\begin{equation}
p_i=a-q_i-\theta q_j,
\qquad i\neq j,
\label{eq:inverse-demand}
\end{equation}
with $0\leq\theta\leq1$. The parameter $\theta$ measures product substitutability and hence the intensity of product-market rivalry. Consumer surplus is
\begin{equation}
CS=\frac12(q_A^2+q_B^2)+\theta q_Aq_B.
\label{eq:cs}
\end{equation}

\subsection{Technology, investment, and emissions}

Firm $i$ has marginal cost
\begin{equation}
c_i=c-x_i-\mu g_i-\nu h_i.
\label{eq:mc}
\end{equation}
Conventional investment $x_i$ raises productivity. Green investment $g_i$ reduces emissions and, when $\mu>0$, also improves competitiveness. Shared green infrastructure $h_i$ is a local public input that can lower firms' marginal cost at rate $\nu>0$. Investment costs are separable and quadratic,
\begin{equation}
C_i(x_i,g_i)
=\frac{k_x}{2}x_i^2+\frac{k_g}{2}g_i^2,
\qquad k_x,k_g>0.
\label{eq:investment-cost}
\end{equation}
There is no exogenous investment budget linking $x_i$ and $g_i$; any interaction between the two investments is generated endogenously through the product market.

Firm $i$'s profit is
\begin{equation}
\pi_i
=(p_i-c_i)q_i
-\frac{k_x}{2}x_i^2
-\frac{k_g}{2}g_i^2
+s_i g_i.
\label{eq:profit}
\end{equation}
Territorial emissions are
\begin{equation}
E_i=e q_i-\beta g_i-\xi h_i,
\qquad e,\beta>0,\quad \xi\geq0.
\label{eq:emissions}
\end{equation}
The term $e q_i$ captures the scale effect of production, while $\beta g_i$ and $\xi h_i$ capture private and public abatement channels.

\subsection{Local-government objective}

The local government values one half of integrated-market consumer surplus, local producer surplus, and progress toward a territorial emissions target. Its objective is
\begin{equation}
W_i
=\frac12CS+PS_i-\frac{\kappa}{2}h_i^2
-\frac d2(E_i-\bar E_i)^2,
\qquad \kappa,d>0,
\label{eq:government-welfare}
\end{equation}
where
\begin{equation}
PS_i=q_i^2-\frac{k_x}{2}x_i^2-\frac{k_g}{2}g_i^2.
\label{eq:producer-surplus}
\end{equation}
At the Cournot stage, operating profit equals $q_i^2$. The subsidy payment $s_i g_i$ is an intrajurisdictional transfer and therefore cancels from regional welfare in the baseline. The quadratic emissions term should be interpreted as a reduced-form loss from deviating from a delegated or territorial decarbonization objective, not as local physical climate damages from carbon emissions.

\subsection{Regularity}

Define
\begin{equation}
R\equiv\frac1{k_x}+\frac{\mu^2}{k_g}.
\label{eq:R}
\end{equation}
We focus on interior equilibria and impose the convenient sufficient condition
\begin{equation}
R<\frac34.
\label{eq:regularity}
\end{equation}
This condition is stronger than necessary, but it guarantees strict concavity of the firms' investment problem and nonsingularity of the reduced output system for all $\theta\in[0,1]$. Government-side strict concavity is imposed where the relevant best-response comparative statics are evaluated. The numerical witness used below satisfies these conditions with slack.
